\documentclass[a4paper,12pt]{article}
\usepackage[utf8]{inputenc}
\usepackage[T1]{fontenc}
\usepackage[ngerman]{babel}
\usepackage{amsmath}
\usepackage{amssymb}
\usepackage{amsthm}
\usepackage{enumitem}
\usepackage{geometry}
\geometry{a4paper, margin=2.5cm}

\title{Einsendeaufgaben -- Lektion 1\\[0.5em]
\large Modul 61112: Lineare Algebra}
\author{}
\date{}

\begin{document}
\maketitle

\section*{Aufgabe 1.3}

\textbf{(UR1)}

\textbf{(UG1)} Da das Nullelement in $K[X]$ liegt, liegt das Nullelement wegen (I2) in $I$.
Insbesondere gilt $R \subseteq R + I$. Da $R \neq \emptyset$, gilt auch $R + I \neq \emptyset$.

\textbf{(UG2)} Sei $x, y \in R + I$ beliebig. Nach der Definition von $R + I$ gilt:
\[
  \exists f_1 \in R\; \exists h_1 \in I \colon x = f_1 + h_1
\]
und
\[
  \exists f_2 \in R\; \exists h_2 \in I \colon y = f_2 + h_2.
\]

Es gilt:
\begin{align*}
  x + y &= f_1 + h_1 + f_2 + h_2 \\
        &= (f_1 + f_2) + (h_1 + h_2)
\end{align*}

Nach (I1) gilt $h_1 + h_2 \in I$ und nach (UG2) gilt $f_1 + f_2 \in R$. Es folgt
\[
  x + y \in R + I
\]

\textbf{(UG3)} Sei $x \in R + I$ beliebig. Nach der Definition von $R + I$ gilt:
\[
  \exists f \in R\; \exists h \in I \colon x = f + h
\]

Das Inverse von $x$ ist $-f - h$.
$-f - h$ liegt in $R + I$, denn $-f \in R$ wegen (UG3) und $(-1) \cdot h \in I$ wegen (I2).

\textbf{(UR2)} Da $1 \in R$ wegen (UR2) und $0 \in I$ wegen (I2): $0 \cdot h = 0$,
folgt $1 + 0 = 1 \in R + I$.

\textbf{(UR3)} Sei $x, y \in R + I$ beliebig. Nach Definition von $R + I$ gilt:
\[
  \exists f_1 \in R\; \exists h_1 \in I \colon x = f_1 + h_1
\]
und
\[
  \exists f_2 \in R\; \exists h_2 \in I \colon y = f_2 + h_2.
\]

Es gilt:
\begin{align*}
  x \cdot y &= (f_1 + h_1) \cdot (f_2 + h_2) \\
            &= f_1 f_2 + f_1 h_2 + f_2 h_1 + h_1 h_2
\end{align*}

Nach (I2) folgt $f_1 h_2 \in I$, $f_2 h_1 \in I$ und $h_1 h_2 \in I$. Wegen (I1) gilt demnach
\[
  f_1 h_2 + f_2 h_1 + h_1 h_2 \in I.
\]

Außerdem gilt $f_1 f_2 \in R$ wegen (UR3). Daraus lässt sich schließen
\[
  x \cdot y \in R + I
\]

\end{document}
